\section{Milvus}

\begin{frame}{Milvus: Distributed Shared-storage Architecture}
	\begin{columns}
		\column{0.5\textwidth}
		\begin{block}{Bốn lớp kiến trúc}
			\begin{itemize}
				\item \textbf{Access Layer}: Proxy stateless, validate và route request.
				\item \textbf{Coordinator}: RootCoord, DataCoord, QueryCoord.
				\item \textbf{Worker Nodes}: DataNode, QueryNode, IndexNode.
				\item \textbf{Storage}: etcd, MinIO/S3, Message Queue.
			\end{itemize}
		\end{block}
		\column{0.5\textwidth}
		\begin{figure}
			\centering
			\fbox{\begin{minipage}{0.9\linewidth}
				\centering
				Proxy\\
				$\downarrow$\\
				RootCoord / DataCoord / QueryCoord\\
				$\downarrow$\\
				DataNode / QueryNode / IndexNode\\
				$\downarrow$\\
				etcd + Object Storage + MQ
			\end{minipage}}
			\caption{Disaggregated architecture của Milvus.}
		\end{figure}
	\end{columns}
	\note{Milvus được thiết kế theo hướng distributed-first. Proxy nhận request, coordinator quản lý metadata và topology, worker nodes xử lý ingestion, search và build index. Storage layer tách riêng giúp hệ thống mở rộng dung lượng và compute độc lập hơn.}
\end{frame}

\begin{frame}{Disaggregated Scalability \& Knowhere Engine}
	\begin{columns}
		\column{0.52\textwidth}
		\begin{block}{Scale theo bottleneck}
			\begin{itemize}
				\item Search chậm $\rightarrow$ thêm QueryNode.
				\item Ingestion chậm $\rightarrow$ thêm DataNode.
				\item Index build chậm $\rightarrow$ thêm IndexNode.
				\item Storage tăng $\rightarrow$ mở rộng MinIO/S3.
			\end{itemize}
		\end{block}
		\column{0.48\textwidth}
		\begin{exampleblock}{Knowhere}
			Abstraction layer cho Faiss, HNSWlib, ScaNN, DiskANN; hỗ trợ SIMD path và GPU acceleration khi có CUDA.
		\end{exampleblock}
		\begin{alertblock}{So với monolithic}
			Qdrant/Weaviate đơn giản hơn; Milvus scale tốt hơn khi corpus vượt hàng trăm triệu vectors.
		\end{alertblock}
	\end{columns}
	\note{Disaggregated architecture cho phép xử lý đúng điểm nghẽn thay vì scale toàn bộ hệ thống. Knowhere làm Milvus linh hoạt về execution engine và index type. Đổi lại là hệ thống có nhiều service và cần năng lực vận hành cao hơn.}
\end{frame}

\begin{frame}{Data Lifecycle \& Operational Gotchas}
	\begin{block}{Lifecycle chuẩn}
		\centering
		\texttt{insert()} $\rightarrow$ Growing Segment $\rightarrow$ \texttt{flush()} $\rightarrow$ Sealed Segment $\rightarrow$ build index $\rightarrow$ \texttt{load()} $\rightarrow$ \texttt{search()}
	\end{block}
	\begin{columns}
		\column{0.52\textwidth}
		\begin{exampleblock}{Growing vs Sealed}
			Growing Segment có thể searchable bằng brute-force; Sealed Segment được persist ra object storage và có full index.
		\end{exampleblock}
		\column{0.48\textwidth}
		\begin{alertblock}{Flush fragmentation}
			Gọi \texttt{flush()} quá thường xuyên với batch nhỏ tạo nhiều sealed segments, gây I/O amplification và giảm throughput.
		\end{alertblock}
	\end{columns}
	\note{Milvus có lifecycle rõ ràng hơn hai hệ thống còn lại. Insert thường là asynchronous, flush đóng segment, load đưa index vào QueryNode RAM. Khi benchmark, cần tách riêng insert, flush, load và search để tránh kết luận sai về latency.}
\end{frame}

\begin{frame}{Milvus HA \& Fault Tolerance}
	\begin{columns}
		\column{0.5\textwidth}
		\begin{block}{HA primitives}
			\begin{itemize}
				\item etcd quorum cho metadata consistency.
				\item Coordinator failover qua service discovery.
				\item QueryNode replica cho serving availability.
				\item Message Queue lưu ingestion stream và WAL-like durability.
			\end{itemize}
		\end{block}
		\column{0.5\textwidth}
		\begin{exampleblock}{Advanced features}
			DiskANN cho dataset vượt RAM, Time Travel theo timestamp, CDC cho replication, GPU Index cho build/search acceleration.
		\end{exampleblock}
		\begin{alertblock}{Production cost}
			HA cần monitoring, backup policy, resource planning và version compatibility.
		\end{alertblock}
	\end{columns}
	\note{Milvus phù hợp enterprise RAG hoặc corpus rất lớn, nơi High Availability và scale độc lập quan trọng hơn sự đơn giản. Các thành phần như etcd, Message Queue và replica giúp chịu lỗi, nhưng đồng thời làm tăng chi phí vận hành.}
\end{frame}
